\color{uniblue}\section[Выбор уставок функции контроля исправности цепей напряжения стороны НН трансформатора]{\sloppy Выбор уставок функции контроля исправности цепей напряжения стороны НН трансформатора} 
\color{black}

\begin{enumerate}[label=\arabic{section}.\arabic*, labelsep=4pt, leftmargin=0pt, itemindent=57pt]

\item
Функция содержит ИО максимального напряжения обратной последовательности и минимального напряжения линейных напряжений стороны НН. КЦН НН срабатывает при повышении напряжения обратной последовательности или при снижении линейных напряжений -- определение любого из этих ненормальных режимов через выдержку времени (отстройка от действия резервных защит) косвенно указывает на неисправность в цепях напряжения.
\item
Функция КЦН НН используется для автоматического переключения режимов работы (блокировка ступени или загрубление уставки по току в зависимости от выбранного режима работы) ступеней максимальной токовой защиты трансформатора, работающими с введенным пуском по напряжению.  

\item
В соответствии с \cite{stofsk1} уставку ИО максимального напряжения обратной последовательности можно определить по выражению:
\begin{equation}
    \label{eq:kznnn1}
    U_{\textsl{2ср}} = k_{\textsl{отс}} \cdot U_{\textsl{обр}},
\end{equation}
\begin{eqexpl}[15mm]
    \item{$U_{\textsl{обр}}$} напряжение обратной последовательности при обрыве одного из проводов цепей ТН, В (принимается равным $U_{\textsl{обр}}=19$ В);
    \item{$k_{\textsl{отс}}$} коэффициент отстройки, принимается равным 0,8.
\end{eqexpl}

Принимаем $U_{\textsl{2ср}} = 0,8 \cdot 19 = 15,2$ В.

\item
В соответствии с \cite{stofsk1} уставка ИО минимального линейного напряжения выбирается по условию отстройки от минимального рабочего режима и может быть принята равной 70 В. 

\item
Выдержка времени функции КЦН НН выбирается на ступень селективности больше времени срабатывания максимальной токовой защиты трансформатора:
\begin{equation}
    \label{eq:zp2}
    t_{\textsl{ср}} = t_{\textsl{ср.МТЗ}}+{\Delta t},
    \end{equation}
    \begin{eqexpl}[15mm]
    \item{$I_{\textsl{ср.МТЗ}}$} время срабатывания МТЗ трансформатора, с;
    \item{$\Delta t$} ступень селективности (принимается равной 0,3-0,5), с.
    \end{eqexpl}

\end{enumerate}
    